%!TeX root=MemoriaTFG.tex

\chapter{Materiales y métodos}
\label{cap:metodos}

Este capítulo describe lo necesario para reproducir el
capítulo~\ref{cap:resultados}. Se presenta primero el material: los
doce datasets externos, la ontología que los unifica bajo un
vocabulario común y el dataset propio DermapixelAI~1.0. A
continuación se enumeran los catorce modelos fundacionales
evaluados, las dos líneas de base convolucionales y la
infraestructura de cómputo. El capítulo cierra con los protocolos
experimentales, las métricas, las pruebas estadísticas y las
medidas de trazabilidad que permiten reconstruir cualquier cifra a
partir de los artefactos publicados.

\section{Datasets}
\label{sec:datos}

La selección de datasets persigue medir la capacidad de
generalización de los modelos entre modalidades de imagen, dominios
clínicos y tareas. Por ello se emplean doce datasets externos
públicos que cubren dermatoscopia, fotografía clínica,
histopatología, análisis de equidad y aprendizaje basado en
conceptos, agrupados por modalidad y finalidad. La
tabla~\ref{tab:datasets-resumen} resume sus características.

\begin{table}[H]
\centering
\footnotesize
\setlength{\tabcolsep}{4pt}
\renewcommand{\arraystretch}{0.95}
\caption{Datasets externos empleados en la evaluación.
$N_{\mathrm{test}}$: tamaño nominal del \emph{split} de evaluación
canónico; ``Clases'': número de categorías diagnósticas o conceptos
en la convención propia del dataset. El asterisco
(\textsuperscript{$\ast$}) marca los datasets con solapamiento
documentado frente al dataset de preentrenamiento Derm1M
(auditoría en sección~\ref{sec:auditoria}). Las dos variantes de Derm7pt
(clínica y dermatoscópica) comparten origen
(Seven-Point Checklist) y se reportan como un único dataset en el
conteo agregado de doce datasets externos.}
\label{tab:datasets-resumen}
\begin{tabular}{lrcll}
\hline
Dataset & $N_{\mathrm{test}}$ & Clases & Modalidad & Uso \\
\hline
HAM10000~\cite{ham10000}             & 1\,232  & 7        & Dermoscopia      & LP, FT, ZS, LLMs \\
BCN20000~\cite{bcn20000}             & 1\,242  & 9        & Dermoscopia      & LP \\
PAD-UFES-20~\cite{padufes}           & 461     & 6        & Clínica móvil    & LP, FT \\
DDI~\cite{ddi}                       & 137     & 2        & Clínica          & LP, FT, equidad \\
Dermnet~\cite{dermnet}               & 4\,002  & 23       & Clínica          & LP\textsuperscript{$\ast$} \\
Derm7pt clínico~\cite{kawahara2019}  & 168     & 2        & Clínica          & LP \\
Derm7pt dermo~\cite{kawahara2019}    & 225     & 2        & Dermoscopia      & LP \\
HIBA                                 & 334     & 2        & Dermoscopia      & LP\textsuperscript{$\ast$} \\
MSKCC                                & 1\,664  & 2        & Dermoscopia      & LP, FT\textsuperscript{$\ast$} \\
WSI patches                          & 12\,354 & 16       & Histopatología   & LP \\
ISIC2018~\cite{isic2018}             & 2\,594  & binaria  & Dermoscopia      & Segmentación \\
Fitzpatrick17k~\cite{fitzpatrick17k} & 16\,577 & 114 + FP & Clínica          & Equidad \\
SkinCon~\cite{skincon}               & 3\,230  & 48 conc. & Clínica          & CBM \\
\hline
\end{tabular}
\end{table}

Los datasets cubren las tres modalidades del diagnóstico
dermatológico. La dermatoscopia (HAM10000, BCN20000, Derm7pt dermo,
HIBA, MSKCC, ISIC2018) aporta imagen polarizada adquirida con
dermatoscopio. La fotografía clínica (PAD-UFES-20, DDI, Dermnet,
Derm7pt clínico) reproduce las condiciones de captura en consulta y
atención primaria. La histopatología la aporta WSI~patches, con
$12\,354$ \emph{parches} en $16$ categorías. Dos datasets cumplen
un papel específico: Fitzpatrick17k habilita el análisis de equidad,
al anotar conjuntamente patología ($114$ categorías) y fototipo
cutáneo (escala Fitzpatrick I-VI), y SkinCon aporta anotación densa
de $48$ conceptos clínicos visuales para los Concept Bottleneck
Models. En todos los casos se respetan las particiones
\emph{train/val/test} canónicas de cada dataset, sin reasignación.

\section{Ontología jerárquica L1/L2/L3}
\label{sec:ontologia}

Comparar modelos entre datasets exige resolver un problema previo:
cada conjunto usa su propia taxonomía, y esas taxonomías no son
compatibles entre sí. Dos diagnósticos equivalentes pueden aparecer
bajo denominaciones distintas o con niveles de granularidad
diferentes. Para unificarlas se diseñó, en colaboración
con la Dra.~R.~Taberner, una ontología jerárquica de tres niveles
que actúa como vocabulario común a todos los protocolos del trabajo.

El nivel L1 (etiológico) reparte los diagnósticos en cuatro
familias: patología inflamatoria, tumoral, infecciosa y
genodermatosis. El nivel L2 (subcategoría clínica) contiene
$43$ entradas; cinco de ellas no se desagregan en L3 por
considerarse ya suficientemente específicas (\emph{Esclerosis
tuberosa}, \emph{Neurofibromatosis tipo~1}, \emph{Incontinentia
pigmenti}, \emph{Queratodermia palmoplantar hereditaria} y
\emph{Pseudoxantoma elástico}). El nivel L3 (diagnóstico) contiene
$367$ entradas individuales, cada una con sus códigos SNOMED~CT y
CIE-10 para garantizar interoperabilidad con sistemas clínicos.

Once de los doce datasets se mapean a la ontología; SkinCon queda
fuera por estar anotado con conceptos, no con diagnósticos. El
conjunto armonizado resultante reúne $72\,654$ imágenes etiquetadas
sobre vocabulario común.

\section{Dataset DermapixelAI 1.0}
\label{sec:dermapixel}

DermapixelAI~1.0 es un dataset dermatológico en español construido
en este trabajo a partir del archivo del blog Dermapixel de la
Dra.~R.~Taberner; su construcción materializa el objetivo~O2
(sección~\ref{sec:objetivos}). A diferencia de los datasets públicos
utilizados en la evaluación, DermapixelAI~1.0 incorpora texto
narrativo clínico en español asociado a cada caso, lo que permite
estudiar escenarios multimodales y de transferencia lingüística no
cubiertos por los conjuntos de referencia internacionales. La
elaboración comprende la
extracción del material del archivo, la depuración, la
estructuración que vincula para cada caso imagen y modalidad,
metadatos clínicos y texto narrativo, y el etiquetado contra la
ontología jerárquica L1/L2/L3; el pipeline completo se documenta en
el Anexo~\ref{cap:anexoc}. La versión utilizada reúne $1\,089$
imágenes vinculadas a $669$ casos clínicos, cada uno con su texto
narrativo asociado (mediana de $223$ palabras). Predomina la
fotografía clínica ($1\,036$ imágenes), con un subconjunto
dermatoscópico reducido ($49$).

Las imágenes se etiquetan según la ontología L1/L2/L3 y se
particionan bajo regla \emph{case-aware} ($891 / 157 / 41$
imágenes; $538 / 100 / 31$ casos), de modo que todas las imágenes
de un mismo caso caen en el mismo \emph{split} y no hay fuga entre
conjuntos. Dos campos cuantifican la calidad del etiquetado: el
$97{,}89\%$ del dataset está mapeado a la ontología con revisión
experta del vocabulario (\texttt{label\_source = ontology}) y el
$84{,}39\%$ ha pasado además revisión visual caso-a-caso por la
Dra.~R.~Taberner (\texttt{rosa\_verified}). Ambas cifras no son
intercambiables y se detallan, junto a un tercer campo de validación
textual, en el Anexo~\ref{cap:anexoc}.

El dataset se distribuye bajo licencia CC-BY-NC-SA~4.0. Su pipeline
de construcción, caracterización completa y datasheet figuran en el
Anexo~\ref{cap:anexoc}, y el documento formal de entrega, en el
Anexo~\ref{cap:anexoe}. La validación experimental sobre
DermapixelAI~1.0 forma parte de los resultados de este trabajo: las
secciones~\ref{sec:dermapixel-results} y~\ref{sec:dermapixel-zs} del
capítulo~\ref{cap:resultados} evalúan sobre este dataset el sondeo
lineal supervisado y la clasificación zero-shot multimodal en los
tres niveles ontológicos.

\section{Modelos evaluados}
\label{sec:modelos}

Los modelos incluidos en el estudio se seleccionan con el objetivo
de representar los principales paradigmas actualmente utilizados en
visión artificial y aprendizaje multimodal aplicado a dermatología.
La comparación combina modelos específicamente entrenados para
dermatología, modelos visuales generalistas de referencia,
arquitecturas visión-lenguaje biomédicas y modelos multimodales de
propósito general. Esta diversidad permite evaluar hasta qué punto
la especialización dermatológica aporta ventajas medibles frente a
modelos más generales y situar el rendimiento de los modelos
fundacionales dermatológicos dentro del ecosistema actual de
inteligencia artificial médica.

La tabla~\ref{tab:modelos} resume los dieciséis sistemas evaluados:
catorce modelos fundacionales ---encoders visuales de dominio
(PanDerm) y generales (DINOv2), modelos visión-lenguaje (la familia
DermLIP, BiomedCLIP y SigLIP), el segmentador promptable SAM2.1 y
los modelos de lenguaje multimodales (GPT-4o, Gemini~2.5~Pro,
MedGemma y BLIP-2)--- y dos líneas de base de redes convolucionales
supervisadas (ConvNeXt-Large y EfficientNetV2-Large), incluidas como
referencia no fundacional frente a la que medir el valor del
preentrenamiento.

\begin{table}[H]
\centering
\footnotesize
\setlength{\tabcolsep}{4pt}
\renewcommand{\arraystretch}{0.95}
\caption{Modelos evaluados.}
\label{tab:modelos}
\begin{tabular}{llrl}
\hline
Modelo & Familia & Parámetros & Referencia \\
\hline
PanDerm Base                & Encoder visual derm.       & $\sim 86$\,M & \cite{panderm2025} \\
PanDerm Large               & Encoder visual derm.       & $\sim 307$\,M & \cite{panderm2025} \\
DINOv2 ViT-L/14             & Encoder visual general     & $\sim 300$\,M & \cite{dinov2} \\
ConvNeXt-Large              & ConvNet supervisada        & $\sim 198$\,M & --- \\
EfficientNetV2-Large        & ConvNet supervisada        & $\sim 118$\,M & --- \\
DermLIP v1 (PubMedBERT)     & Vision-lenguaje derm.      & --- & \cite{derm1m} \\
DermLIP v2 (PubMedBERT)     & Vision-lenguaje derm.      & --- & \cite{derm1m} \\
DermLIP original (GPT-2)    & Vision-lenguaje derm.      & --- & \cite{derm1m} \\
BiomedCLIP                  & Vision-lenguaje biomédico  & --- & \cite{biomedclip} \\
SigLIP-Large SO400M         & Vision-lenguaje general    & $\sim 878$\,M & \cite{siglip2023} \\
SAM2.1-Large (decoder FT)   & Segmentación               & $\sim 227$\,M & \cite{sam2} \\
GPT-4o                      & LLM multimodal             & no público & --- \\
Gemini 2.5 Pro              & LLM multimodal             & no público & --- \\
MedGemma 4B Vision          & LLM multimodal biomédico   & $\sim 4$\,B & --- \\
MedGemma 27B Text+LoRA      & LLM multimodal biomédico   & $\sim 27$\,B & --- \\
BLIP-2 (Q-Former)           & Vision-lenguaje general    & --- & \cite{blip2_2023} \\
\hline
\end{tabular}
\end{table}

Los encoders difieren en la dimensión de sus \emph{embeddings}:
$768$ en PanDerm Base; $1\,024$ en PanDerm Large y DINOv2
ViT-L/14; $1\,152$ en SigLIP-Large SO400M. La cifra de
$\sim 307$\,M parámetros de PanDerm Large es el valor nominal del
ViT-L/16 publicado por sus autores. Al cargarlo como extractor de
características y sustituir la cabeza original por la identidad
quedan $303{,}3$\,M parámetros efectivos. Esta diferencia no afecta
al rendimiento, porque la cabeza descartada no interviene en la
inferencia. La familia DermLIP usa ViT-B/16 como encoder visual y
dos encoders textuales alternativos: PubMedBERT extendido a $256$
tokens (variantes v1 y v2) y GPT-2 con tokenizador de CLIP a $77$
tokens (variante original). DermFM-Zero~\cite{dermfmzero2026} se
incluye en la discusión pero no en la evaluación empírica, ya que
sus pesos no estaban disponibles a la fecha de cierre.

\section{Infraestructura experimental}
\label{sec:infraestructura}

Todos los experimentos corren sobre una NVIDIA DGX Spark con chip
Grace Hopper GB10, $128$\,GB de memoria unificada compartida entre
CPU y GPU, Linux ARM aarch64 y CUDA~13.0. La precisión nativa es
BF16 (\emph{brain floating-point} de 16 bits) en inferencia y
entrenamiento, con FP32 en optimizador y pérdida durante el ajuste
fino (precisión mixta canónica). La memoria unificada es clave para
este trabajo: permite mantener cargados modelos del orden de
$55$\,GB en BF16 sin \emph{offloading} ni cuantización, de modo
que incluso MedGemma~27B opera localmente y sin comprimir. El
software queda fijado en Python~3.12.3, PyTorch~2.11.0,
\texttt{timm}~0.9.16, \texttt{scikit-learn}, \texttt{transformers},
\texttt{open\_clip}, \texttt{sam2}, \texttt{peft} y
FAISS~\cite{faiss2019}.

Los tiempos de cálculo característicos sobre este equipo orientan
sobre el coste de cada protocolo: un sondeo lineal completo de un
encoder sobre un dataset estándar tarda $1$--$3$ minutos; el
ajuste fino de PanDerm Large sobre HAM10000 ($50$ épocas con TTA),
unas $15$ horas; el ajuste fino del decodificador de SAM2.1-Large
sobre ISIC2018 ($50$ épocas), unas $24$ horas; la validación
cruzada de cinco pliegues del modelo de segmentación, con la
ablación de rango \emph{LoRA} asociada, unas $18$ horas; y el
entrenamiento de un Sparse Autoencoder Large ($16\,384$
características), unas $4$ horas.

\section{Protocolos experimentales}
\label{sec:protocolos}

Los protocolos experimentales se diseñan para caracterizar
distintas capacidades de los modelos fundacionales dermatológicos
bajo condiciones complementarias. Cada protocolo responde a una
pregunta específica: el sondeo lineal evalúa la calidad intrínseca
de las representaciones aprendidas durante el preentrenamiento; el
ajuste fino cuantifica el rendimiento máximo alcanzable tras
adaptación supervisada; la eficiencia en etiquetas analiza la
cantidad de datos anotados necesaria para obtener resultados
competitivos; la segmentación estudia la localización espacial de
lesiones; la clasificación zero-shot mide la capacidad de
transferencia sin entrenamiento específico; la comparación con
modelos multimodales generalistas sitúa el rendimiento de los
modelos especializados en el contexto actual de los grandes modelos
fundacionales; la evaluación de equidad analiza el comportamiento
sobre diferentes fototipos cutáneos; y la auditoría de solapamiento
con Derm1M permite interpretar correctamente los resultados de
transferencia y zero-shot. En conjunto, estos experimentos
proporcionan una caracterización multidimensional de los modelos
evaluados, más allá de una única métrica o tarea aislada. Todos los
protocolos comparten la infraestructura
(sección~\ref{sec:infraestructura}) y las semillas declaradas
(sección~\ref{sec:reproducibilidad}), de modo que sus cifras son
comparables entre sí.

\subsection{Sondeo lineal (LP)}
\label{sec:lp}

El sondeo lineal mide el valor de las representaciones de un encoder
sin modificarlo: se congela el encoder y se entrena un clasificador
lineal sobre los \emph{embeddings} que extrae. La extracción usa la
normalización canónica de cada modelo —parámetros ImageNet
($\mu=[0{,}485, 0{,}456, 0{,}406]$,
$\sigma=[0{,}229, 0{,}224, 0{,}225]$) para PanDerm y DINOv2, y
parámetros OpenAI CLIP para DermLIP, BiomedCLIP y SigLIP— y los
vectores se guardan en disco para no recomputarlos entre
experimentos. El clasificador es una regresión logística multinomial
(L-BFGS, penalización L2 con $C = 1{,}0$, \texttt{max\_iter}~$=5\,000$,
\texttt{random\_state}~$=42$) entrenada sobre las particiones
canónicas de la tabla~\ref{tab:datasets-resumen}. El protocolo se
aplica a PanDerm Base y Large, DINOv2, ConvNeXt-Large,
EfficientNetV2-Large y la rama visual de DermLIP v2, BiomedCLIP y
SigLIP-Large.

\subsection{Ajuste fino supervisado (FT)}
\label{sec:ft}

A diferencia del sondeo lineal, el ajuste fino sí modifica el
encoder: entrena de forma conjunta el encoder visual y una cabeza de
clasificación. La receta reproduce la del paper original de
PanDerm~\cite{panderm2025}, para que la comparación sea directa.
El optimizador es AdamW (\emph{weight decay} $0{,}05$) con tasa de
aprendizaje inicial $\eta = 5\times 10^{-4}$, \emph{cosine
annealing} y \emph{warmup} lineal de $10$ épocas. La regularización
combina \emph{layer decay} $0{,}65$, \emph{drop-path} $0{,}2$,
\emph{label smoothing} $\epsilon = 0{,}1$, Mixup $\alpha = 0{,}8$
y CutMix $\alpha = 1{,}0$, junto a un \emph{weighted random
sampler} con pesos inversos a la frecuencia de clase. El
entrenamiento dura $50$ épocas en precisión mixta canónica, con
las semillas de \texttt{torch}, \texttt{numpy} y \texttt{random}
fijadas a $0$. En la evaluación final sobre HAM10000 se añade
\emph{test-time augmentation} (TTA): se promedian los logits de
$5$ augmentaciones deterministas por imagen (volteos horizontal y
vertical, rotación y \emph{jitter} de color). El protocolo se
aplica a PanDerm Base y Large sobre HAM10000, PAD-UFES-20, DDI y
MSKCC.

\subsection{Eficiencia en etiquetas}
\label{sec:label-eff}

Este protocolo responde a una pregunta práctica: cuántas etiquetas
hacen falta. Manteniendo el encoder congelado, se entrenan
clasificadores lineales (misma configuración que el sondeo lineal,
sección~\ref{sec:lp}) sobre submuestreos estratificados del
conjunto de entrenamiento de HAM10000 y PAD-UFES-20 al $1\%$,
$5\%$, $10\%$, $20\%$, $50\%$ y $100\%$. La
estratificación preserva la proporción de clases en cada subconjunto
y los submuestreos se generan con \texttt{random\_state}~$=42$.
PanDerm Large es el encoder principal; PanDerm Base se usa como
referencia adicional sobre HAM10000.

\subsection{Segmentación binaria de lesión}
\label{sec:segmentacion}

La segmentación constituye una tarea complementaria a la
clasificación, ya que permite evaluar no solo si un modelo reconoce
correctamente una lesión, sino también si es capaz de localizarla
con precisión dentro de la imagen. En el contexto de los modelos
fundacionales recientes basados en la familia \emph{Segment
Anything}, surge además una cuestión relevante: determinar hasta qué
punto el rendimiento observado depende de la arquitectura utilizada,
del preentrenamiento específico en imagen médica o del tamaño del
\emph{backbone} visual. Para responder a estas preguntas, la
evaluación se organiza en dos protocolos complementarios y varios
análisis de robustez.

\paragraph{Adaptación de SAM2.1-Large al dominio dermatológico.}
El primer protocolo adapta SAM2.1-Large mediante \emph{fine-tuning}
denso de su decodificador: se ajustan el
decodificador de máscaras (\texttt{mask\_decoder}) y el
\emph{promptable encoder}, mientras que el \texttt{image\_encoder}
(Hiera-Large preentrenado) permanece congelado, de modo que solo
$\sim 4{,}2$\,M parámetros resultan entrenables ($\sim 17$\,MB de
\emph{checkpoint}). La optimización usa AdamW
($\eta = 1\times 10^{-4}$, \emph{weight decay} $0{,}05$) y una
pérdida que combina \emph{cross-entropy} y \emph{Dice loss} sobre la
máscara binaria; el \emph{batch} es de $4$ (limitado por el coste de
SAM2.1-Large en BF16) y la normalización es uniforme
($\mu = \sigma = 0{,}5$ en los tres canales), según la convención
de SAM2. Se contrastan $50$ épocas (resultado de referencia) y
$100$ con \emph{early stopping} sobre el Dice de validación. El
entrenamiento es sobre ISIC2018~\cite{isic2018} y la transferencia
fuera de dominio se evalúa sobre ISIC2017 y PH2 sin reentrenamiento.
Como referencias no adaptadas se incluyen SAM2 \emph{zero-shot} sin
\emph{fine-tuning} y un \emph{Convolutional AutoEncoder} (CAE-seg)
entrenado desde cero con idéntica pérdida y optimización.

\paragraph{Comparación controlada entre arquitecturas de
segmentación.}
El segundo protocolo separa los tres factores que determinan el
rendimiento ---generación de arquitectura, preentrenamiento de
dominio y tamaño de \emph{backbone}--- mediante una comparativa
pareada de cinco arquitecturas de segmentación \emph{promptable}:
MedSAM2-tiny~\cite{medsam2}, SAM2.1-tiny y SAM2.1-Large (familia
SAM2~\cite{sam2}), y SAM-Med2D~\cite{sammed2d} y MedSAM
v1~\cite{medsam} (familia SAM~\cite{sam2023}). Todas se entrenan
bajo un régimen común (\emph{fine-tuning} denso del decodificador y
el \emph{promptable encoder}, encoder visual congelado, \emph{prompt}
de \emph{bounding box} derivada de la máscara de referencia, AdamW
$\eta = 1\times 10^{-4}$, pérdida Dice$+$BCE, $30$ épocas con
\emph{early stopping}) sobre un dataset dermatoscópico ampliado y
deduplicado por hash MD5 de píxel ($14\,201$ imágenes únicas de
ISIC2018, ISIC2017, HAM10000-seg y PH2), y se evalúan sobre el mismo
test canónico de ISIC2018 ($N = 1\,000$, blindado frente al
entrenamiento). La significación de las diferencias de Dice se
contrasta con el test de Wilcoxon pareado~\cite{wilcoxon1945}.

\paragraph{Análisis complementarios.}
Para interpretar correctamente las diferencias observadas entre
arquitecturas, se realizan además tres análisis complementarios
dirigidos a evaluar la contribución del \emph{prompt}, la
estabilidad del entrenamiento y la robustez de las predicciones.
Para aislar la contribución del \emph{prompt}, un
control automático sin caja ---una U-Net con codificador ResNet-101
preentrenado en ImageNet ($51{,}5$\,M parámetros)--- se entrena
sobre el mismo dataset y se evalúa sin recibir \emph{bounding box}.
La estabilidad de la mejor configuración se estima mediante
validación cruzada de cinco pliegues estratificados por dataset de
origen y una ablación del rango de adaptación de bajo rango
(\emph{LoRA}, $r\in\{4,8,16,32\}$); su robustez frente a la calidad
del \emph{prompt} se mide bajo perturbaciones controladas de la
\emph{bounding box} (desplazamiento de esquinas y cambio de escala),
y la confianza de las máscaras se calibra \emph{a posteriori}
mediante \emph{temperature scaling}. Por último, el techo de la
métrica se contextualiza frente al acuerdo inter-anotador del propio
dataset.

\subsection{Clasificación zero-shot multimodal}
\label{sec:zs}

La clasificación zero-shot constituye una de las capacidades más
características de los modelos fundacionales multimodales. A
diferencia de los enfoques supervisados, estos modelos pueden
realizar tareas de clasificación sin haber sido entrenados
específicamente sobre las clases objetivo, utilizando únicamente
descripciones textuales de las categorías candidatas. Evaluar esta
capacidad permite estimar el grado de generalización adquirido
durante el preentrenamiento y determinar hasta qué punto el
conocimiento dermatológico puede transferirse a nuevos conjuntos de
datos sin adaptación adicional.

El zero-shot mide qué pueden clasificar los modelos vision-lenguaje
sin ningún entrenamiento específico, aprovechando su espacio
compartido imagen-texto. Cada clase candidata se formula como una o
varias frases de plantilla, que el encoder textual codifica y
—cuando hay varias por clase— promedia en un único vector. La
imagen se codifica con el encoder visual y gana la clase de mayor
similitud coseno. Sobre HAM10000 se comparan tres formulaciones de
\emph{prompt}: (i)~genéricos en inglés, una plantilla por clase
(p.~ej.~\emph{``a dermoscopic image of a melanoma''});
(ii)~el conjunto A3 del paper de Derm1M~\cite{derm1m}, con varias
plantillas por clase; y (iii)~una versión enriquecida de A3 con
descriptores clínicos (modalidad, localización, criterios visuales).
El conjunto A3 es la estrategia de \emph{prompting} propuesta en
Derm1M~\cite{derm1m} que sustituye el nombre simple de cada
enfermedad por un conjunto de descripciones clínicas y visuales más
completas: en lugar de representar una clase únicamente con la
palabra \emph{melanoma}, emplea varias frases que describen sus
características diagnósticas, lo que produce representaciones
textuales más informativas y mejora el rendimiento \emph{zero-shot}.
Para caracterizar la sensibilidad al tokenizador, el protocolo se
ejecuta con los dos tokenizadores de la familia DermLIP: PubMedBERT
a $256$ tokens (v1/v2) y GPT-2 con tokenizador de CLIP a $77$
tokens (original). Sobre HIBA, MSKCC, DDI y Derm7pt clínico se corre
la versión binaria melanoma/no-melanoma con \emph{prompts}
genéricos y A3. De este modo, el protocolo permite aislar el efecto
de tres componentes fundamentales del rendimiento zero-shot: el
modelo multimodal utilizado, la formulación textual de las clases y
el tokenizador empleado por la rama lingüística.

\subsection{Comparación con modelos de lenguaje multimodales}
\label{sec:llm}

La rápida evolución de los modelos multimodales de propósito general
plantea una cuestión relevante para la dermatología computacional:
determinar si los grandes modelos fundacionales capaces de procesar
simultáneamente imagen y texto pueden competir con modelos
especializados entrenados específicamente para imagen dermatológica.
Este protocolo evalúa dicha hipótesis comparando varios modelos
multimodales de última generación frente a los encoders
dermatológicos analizados en las secciones anteriores bajo un
escenario común de clasificación diagnóstica. La comparación
contrapone así dos paradigmas ---la especialización dermatológica
explícita de los encoders de las secciones anteriores frente a la
escala de los grandes modelos generalistas---, con los modelos
biomédicos generales (MedGemma) en una posición intermedia, lo que
permite contrastar la contribución relativa de la especialización de
dominio frente a la escala del preentrenamiento.

Para situar a los encoders frente a los grandes modelos
generalistas, el protocolo evalúa GPT-4o (API OpenAI),
Gemini~2.5~Pro (API Google), MedGemma~4B~Vision (local) y
MedGemma~27B Text (local en BF16) sobre el mismo \emph{split} de
test de HAM10000 ($1\,232$ imágenes) usado en sondeo lineal y
ajuste fino. El \emph{prompt} clínico estandarizado (unas $300$
palabras) plantea la tarea como apoyo al diagnóstico —no
sustitución del juicio clínico—, presenta las siete categorías de
HAM10000 con una descripción breve y pide el diagnóstico más
probable con razonamiento estructurado. La etiqueta se extrae por
\emph{parsing} determinista que mapea el nombre del diagnóstico al
código de HAM10000; las respuestas no mapeables (menos del $1\%$
en todos los modelos) cuentan como incorrectas. El modo principal es
\emph{zero-shot}; para MedGemma~27B se evalúa además una variante
ajustada con LoRA ($r = 8$) sobre el entrenamiento de HAM10000.

\subsection{Evaluación de equidad por fototipo}
\label{sec:equidad}

La equidad constituye una dimensión relevante en la evaluación de
sistemas de inteligencia artificial aplicados a dermatología.
Diversos estudios han señalado que muchos conjuntos de entrenamiento
presentan una representación desigual de los distintos fototipos
cutáneos, lo que puede traducirse en diferencias de rendimiento
entre grupos poblacionales. En consecuencia, además de evaluar la
precisión global de los modelos, resulta necesario analizar si su
comportamiento se mantiene estable a lo largo de todo el espectro de
fototipos definido por la escala de Fitzpatrick.

Este protocolo mide si el rendimiento se mantiene entre fototipos
cutáneos. Cinco encoders (PanDerm Large, PanDerm Base, la rama
visual de DermLIP v2, DINOv2 ViT-L/14 y BiomedCLIP) se evalúan sobre
Fitzpatrick17k completo ($16\,577$ imágenes, $114$ patologías,
$6$ fototipos I-VI) por sondeo lineal (configuración de
sección~\ref{sec:lp}, con la normalización canónica de cada
modelo). Se reportan dos formulaciones: clasificación a $114$
patologías (grano fino) y a $3$ clases de malignidad
(benigno / no neoplásico / maligno), esta última habitual en
estudios de equidad sobre Fitzpatrick17k. Las métricas se calculan
por subgrupo sobre los seis fototipos, y la equidad se resume con la
brecha entre los extremos
$\mathrm{Gap}_{\mathrm{I}\to\mathrm{VI}} =
\mathrm{BAcc}_{\mathrm{I}} - \mathrm{BAcc}_{\mathrm{VI}}$. La
brecha entre fototipos extremos resume de forma sencilla la
estabilidad del rendimiento entre grupos; su interpretación, no
obstante, debe realizarse conjuntamente con la exactitud global, ya
que una diferencia reducida puede reflejar tanto un comportamiento
equitativo como un rendimiento uniformemente bajo en todos los
subgrupos. El experimento se verifica re-ejecutando el sondeo sobre
PanDerm Large y comprobando coincidencia hasta cuatro decimales con
el resultado original.

\subsection{Auditoría de solapamiento con Derm1M}
\label{sec:auditoria}

La interpretación de los resultados zero-shot requiere verificar
previamente la independencia entre los conjuntos de evaluación y los
datos utilizados durante el preentrenamiento. Si un modelo ha visto
durante su entrenamiento imágenes posteriormente utilizadas como
test, las métricas obtenidas pueden reflejar memorización parcial
del conjunto de evaluación y no capacidad real de generalización.
Por este motivo, se realiza una auditoría sistemática de
solapamiento entre los datasets evaluados y Derm1M, principal
dataset de preentrenamiento de la familia PanDerm y DermLIP.

La auditoría lo cuantifica: cada imagen se reduce
a su hash MD5 (biblioteca estándar de Python) y la intersección de
hashes identifica las imágenes presentes en ambos datasets. La
cardinalidad se reporta por dataset en el
capítulo~\ref{cap:resultados}. Se auditan los doce datasets
externos, salvo SkinCon (no evaluado sobre clases diagnósticas).

El alcance de esta auditoría se declara con precisión. El hash MD5
detecta coincidencias \emph{exactas} a nivel de bytes y, por tanto,
identifica duplicados idénticos, pero no \emph{near-duplicates}
---recortes, reescalados, recompresiones o alteraciones leves de una
misma imagen---. La auditoría sostiene, en consecuencia, la ausencia
de \emph{solapamiento exacto}, no la ausencia completa de
contaminación visual. Esta elección prioriza la trazabilidad y el bajo
coste del hash exacto frente a la cobertura perceptual; el
\emph{perceptual hashing} o la comparación por \emph{embeddings}
visuales constituyen su extensión natural, y las implicaciones de este
alcance acotado para la interpretación de los resultados se retoman en
la sección~\ref{sec:lim-leakage}.

\section{Métricas}
\label{sec:metricas}

Ninguna métrica individual captura por sí sola todos los aspectos
relevantes del rendimiento de un modelo de inteligencia artificial.
En problemas dermatológicos, caracterizados por distribuciones de
clases desbalanceadas, diferentes niveles de granularidad
diagnóstica y requisitos clínicos diversos, resulta necesario
combinar métricas complementarias que evalúen discriminación,
robustez, concordancia, equidad y viabilidad operativa. Por este
motivo, el presente trabajo reporta un conjunto de métricas
específicas para clasificación, segmentación, equidad y rendimiento
computacional.

\paragraph{Clasificación.}
La métrica principal es la \emph{Balanced Accuracy} (BAcc, promedio
del \emph{recall} por clase): se adopta porque asigna el mismo peso
a todas las clases con independencia de su frecuencia, lo que evita
que las categorías mayoritarias dominen la evaluación en los
datasets desbalanceados de este trabajo. La \emph{accuracy} acompaña
como referencia comparativa
entre modelos sobre la misma distribución. Se reportan además la F1
ponderada (media armónica de precisión y \emph{recall}, pesada por
frecuencia de clase), la AUROC \emph{one-vs-rest}
(\texttt{multi\_class='ovr'} en \texttt{scikit-learn}, promediada
sobre clases) y el coeficiente kappa de Cohen, que corrige la
concordancia por azar. \emph{Recall@$k$} —presencia de la etiqueta
verdadera entre las $k$ primeras posiciones del ranking— se reporta
específicamente para melanoma sobre HAM10000.

\paragraph{Segmentación.}
La calidad de las máscaras predichas se evalúa mediante el
coeficiente Dice (DSC) y la intersección sobre unión (IoU), dos
métricas estándar que cuantifican el grado de solapamiento entre la
segmentación automática y la máscara de referencia:
$\mathrm{DSC} = 2\,|P\cap G|/(|P|+|G|)$ e
$\mathrm{IoU} = |P\cap G|/|P\cup G|$, con $P$ la máscara predicha y
$G$ la de referencia.

\paragraph{Equidad.}
BAcc por subgrupo Fitzpatrick (I-VI) sobre la formulación de tres
clases de malignidad, y \emph{gap}
$\mathrm{Gap}_{\mathrm{I}\to\mathrm{VI}}
= \mathrm{BAcc}_{\mathrm{I}} - \mathrm{BAcc}_{\mathrm{VI}}$. Un
valor próximo a cero indica rendimiento comparable entre fototipos
extremos; su interpretación conjunta con la BAcc global se discute en
la sección~\ref{sec:equidad}.

\paragraph{Operativas.}
Además del rendimiento predictivo, se evalúa el coste computacional
asociado a cada modelo, que permite valorar la viabilidad de un
despliegue en entornos clínicos con recursos limitados: latencia por
consulta en milisegundos (inferencia más posproceso) y huella de
memoria (VRAM en gigabytes, BF16), medidas sobre la infraestructura
de sección~\ref{sec:infraestructura}.

\paragraph{Agregación entre datasets.}
\label{sec:agregacion-datasets}
Algunas vistas panorámicas del capítulo~\ref{cap:resultados} reportan
la media de una métrica sobre varios datasets externos. Dicha media se
interpreta como \emph{brújula descriptiva} ---resume una tendencia
sobre conjuntos de distinta taxonomía, modalidad y prevalencia---, no
como un marcador conmensurable que defina un ranking absoluto entre
modelos: HAM10000, DDI, PAD-UFES-20 o Dermnet no son objetos
perfectamente comparables. Por ello toda cifra agregada se acompaña
siempre del desglose por dataset y por nivel ontológico, que es la
base sobre la que se sostienen las conclusiones comparativas; la media
entre tareas heterogéneas se usa para describir, no para ordenar.

\section{Pruebas estadísticas}
\label{sec:estadistica}

La comparación de modelos no debe limitarse a la observación de
diferencias en las métricas de rendimiento. Incluso cuando un modelo
obtiene una puntuación superior a otro, resulta necesario determinar
si dicha diferencia refleja una ventaja real o puede explicarse por
la variabilidad inherente de la muestra utilizada para la
evaluación. Esta consideración es especialmente relevante en
dermatología, donde los conjuntos de datos suelen presentar tamaños
moderados, distribuciones de clases desequilibradas y una elevada
heterogeneidad clínica.

Con este objetivo, el presente trabajo incorpora un protocolo
estadístico explícito, en línea con las guías de reporte para
inteligencia artificial en imagen médica (CLAIM~\cite{claim2024} y
TRIPOD+AI~\cite{tripodai2024}), basado en la estimación de
intervalos de confianza mediante \emph{bootstrap}, la comparación
formal de curvas ROC mediante el test de DeLong y el control del
riesgo de falsos positivos en escenarios con múltiples
comparaciones. De este modo, los resultados presentados en el
capítulo~\ref{cap:resultados} pueden interpretarse como evidencia
estadística reproducible y no únicamente como un ordenamiento de
modelos derivado de una realización concreta de los datos.

La incertidumbre de cada métrica principal se cuantifica mediante
\emph{bootstrap} no paramétrico con $1\,000$ remuestreos del
conjunto de evaluación, estratificado por clase, del que se derivan
los intervalos de confianza al $95\%$. No requiere asumir
normalidad ni varianza conocida, lo que es necesario sobre los
conjuntos de test reducidos y desbalanceados de este trabajo, donde
la aproximación asintótica de los estimadores no es fiable.

Las comparaciones de AUROC entre dos modelos sobre el mismo conjunto
emplean el test de DeLong~\cite{delong1988}. Su motivación es
metodológica: ambos modelos se evalúan sobre las mismas imágenes, de
modo que sus curvas ROC están correlacionadas; tratar las dos AUROC
como muestras independientes subestima la varianza del contraste y
sobrestima su significación. DeLong estima esa covarianza de forma
no paramétrica y produce un contraste pareado válido. Se aplica a
los contrastes pareados entre los cinco encoders sobre Fitzpatrick17k
(sección~\ref{sec:equidad-delong}, tabla~\ref{tab:delong-pares}).

Los $p$-valores se reportan como nominales y, de forma
complementaria, bajo corrección de Bonferroni para el número de
comparaciones. Exponer ambos es intencionado: la corrección controla
la tasa de falsos positivos a costa de potencia estadística, y
mostrar el valor nominal junto al corregido permite al lector juzgar
la robustez de cada contraste sin imponer una única política de
corrección. Dada la magnitud de las diferencias observadas (con $p$
varios órdenes por debajo de $0{,}05$), las conclusiones principales
se mantienen bajo cualquier corrección estándar. Los tamaños
muestrales por subgrupo del experimento de equidad se reportan junto
a sus cifras en el capítulo~\ref{cap:resultados}.

El alcance de este protocolo estadístico es, no obstante, heterogéneo
entre bloques experimentales, y conviene declararlo de forma
explícita. El contraste pareado de DeLong y la corrección por
comparaciones múltiples se aplican donde la comparación entre modelos
es el objeto central ---el ranking de encoders en detección de
malignidad (sección~\ref{sec:equidad-delong})---; la estimación por
\emph{bootstrap} acompaña a las métricas principales sobre
DermapixelAI~1.0 y a la validación cruzada de segmentación. Otros
bloques ---el sondeo lineal sobre los datasets externos, la eficiencia
en etiquetas y la equidad por fototipo--- se reportan como
estimaciones puntuales bajo semilla única
(sección~\ref{sec:lim-semilla}). En consecuencia, la fuerza de las
afirmaciones comparativas se modula según el cierre estadístico
disponible: se sostienen como patrones cualitativos robustos cuando la
magnitud del efecto excede holgadamente la variabilidad esperable, y
se enuncian con la cautela correspondiente cuando el contraste formal
queda fuera del alcance del trabajo.

\section{Trazabilidad y reproducibilidad}
\label{sec:reproducibilidad}

La reproducibilidad no constituye un complemento documental, sino
una condición necesaria para la validación independiente de los
resultados. Un experimento cuyos resultados no pueden reconstruirse
de forma consistente resulta difícil de verificar y limita la
confianza en las conclusiones derivadas. El trabajo se diseña, por
tanto, para que cualquier número del capítulo~\ref{cap:resultados}
se reconstruya a partir de artefactos persistentes y no de la
memoria de una ejecución concreta. Tres mecanismos articulan esta
garantía.

El primero es el determinismo de la ejecución. Las semillas
aleatorias se fijan a $0$ para \texttt{torch}, \texttt{numpy} y
\texttt{random} en \emph{fine-tuning} y segmentación, y a $42$ en
\texttt{random\_state} para los protocolos basados en
\texttt{scikit-learn} (sondeo lineal, eficiencia en etiquetas,
equidad). El entorno queda fijado a las versiones de
sección~\ref{sec:infraestructura}, y cualquier desviación se
documenta en el experimento afectado. Por último, se deshabilita la
telemetría externa y se fija explícitamente el dispositivo de
ejecución para evitar variaciones no controladas del entorno
experimental.

El segundo mecanismo es la persistencia de las salidas. Los
resultados numéricos se vuelcan a ficheros CSV o JSON con esquema
estable, lo que permite reconstruir cualquier tabla o figura del
capítulo~\ref{cap:resultados} sin reejecutar nada, y los mapeos de
etiquetas nativas a la ontología L1/L2/L3 se mantienen fijos entre
experimentos. El tercero es la integridad de la cadena de cómputo:
cuando un experimento toma como entrada la salida de otro ---por
ejemplo, los \emph{embeddings} del sondeo lineal alimentan la rama
visual del zero-shot---, la identidad de esa entrada se preserva
mediante hash MD5 del fichero. El encadenamiento por hash convierte
el conjunto de experimentos en un grafo de dependencias verificable:
cualquier alteración silenciosa de una entrada intermedia
invalidaría el hash y sería detectable, lo que protege la cadena de
evidencia frente a la deriva de datos entre etapas dependientes. El
listado completo de tareas reproducibles, con sus comandos, ficheros
de salida y \emph{checkpoints}, figura en el
Anexo~\ref{cap:anexoa}.
